\section{BT1083--BT1085 Matter Bridge Update}
\label{sec:bt1083-bt1085-matter-bridge}

The newest matter-sector bridge replaces the earlier diagonal rank template by a linear transvection projector.  The BT876 long-root transvection has thirteen fixed basis vectors and nine shell $3$-cycles.  Its grade-zero projector on $\mathbb C[40]$ is
\[
P_{22}=P_{\mathrm{fixed}}+\sum_{C}\frac{1}{3}\mathbf 1_C\mathbf 1_C^{T},
\]
where the sum runs over the nine shell cycles.  Thus
\[
\operatorname{rank}P_{22}=13+9=22,
\qquad
P_{22}^2=P_{22}.
\]
This is a Fourier/eigenspace projector, not a selector of twenty-two points.

The first lift into the $162$-slot matter carrier uses one $54$-dimensional block per generation,
\[
\mathbb C^{162}=G_0\oplus G_1\oplus G_2,
\qquad
G_g=\mathbb C^{40}_{\mathrm{core}}(g)\oplus R_{14}(g).
\]
On each generation block the lifted complement projector is
\[
P_{22}^{(g)}=P_{22}\oplus 0_{14},
\]
so the physical block is
\[
P_{32}^{(g)}=(I_{40}-P_{22})\oplus I_{14},
\qquad
\operatorname{rank}P_{32}^{(g)}=18+14=32.
\]
Consequently
\[
\operatorname{rank}P_{66}=66,
\qquad
\operatorname{rank}P_{96}=96.
\]
The corresponding chain-side rank-$96$ support remains the sector projector
\[
P_{96}^{\mathrm{chain}}=P_0+P_{16},
\qquad
81+15=96.
\]
A sparse bridge skeleton is a partial isometry
\[
F:\mathbb C^{162}\longrightarrow\mathbb C^{240}
\]
which maps the slot-side $96$ block onto $E_0\oplus E_{16}$ and embeds the $66$ complement into $E_4\oplus E_{10}$.  The unused reservoir has dimension
\[
(120+24)-66=78=66+12=6\Phi_3=2(3\Phi_3).
\]
The preferred reading is that the reservoir carries one additional $3\times22$ transvection-bookkeeping scale together with the $1+3+8$ gauge-adjoint profile.  This closes the rank-level matter bridge while leaving the natural incidence/gauge action on the reservoir as the next construction.

For the scalar/gauge ladder, the W33-native line-edge adjacency on $C_1$ gives projected nearest-sector ranks
\[
0\to4:31,\qquad 4\to10:24,\qquad 10\to16:15,
\]
whereas the triangle-edge adjacency gives $38,22,15$.  Thus the line-edge adjacency is the current best ladder candidate: it is full rank on the last two steps and has the smaller oriented gap-square ledger
\[
2\left(31+24\left(\frac32\right)^2+15\left(\frac32\right)^2\right)=\frac{475}{2}.
\]
